\documentclass[12pt,a4paper]{article}
\usepackage{multirow}
\usepackage[utf8]{inputenc}
\usepackage{geometry}
\geometry{margin=0.75in}
\usepackage{mathptmx} % Times New Roman font
\usepackage{helvet}
\usepackage{courier}
\usepackage{amsmath}
\usepackage{amsfonts}
\usepackage{amssymb}
\usepackage{graphicx}
\usepackage{booktabs}
\usepackage{array}
\usepackage{longtable}
\usepackage{url}
\usepackage{xcolor}
\usepackage{hyperref}
\usepackage{subcaption}
\usepackage{float}
\usepackage{algorithm}
\usepackage{algorithmic}
\usepackage{listings}
\usepackage{tcolorbox}
\tcbuselibrary{most}

% Professional color system
\definecolor{primaryblue}{RGB}{31,81,138}
\definecolor{secondaryblue}{RGB}{70,130,180}
\definecolor{accentgreen}{RGB}{40,167,69}
\definecolor{warningred}{RGB}{220,53,69}
\definecolor{highlightgray}{RGB}{248,249,250}
\definecolor{tableheader}{RGB}{233,236,239}

% Professional highlighting commands
\newcommand{\primarytext}[1]{\textcolor{primaryblue}{\textbf{#1}}}
\newcommand{\secondarytext}[1]{\textcolor{secondaryblue}{\textbf{#1}}}
\newcommand{\success}[1]{\textcolor{accentgreen}{\textbf{#1}}}
\newcommand{\warning}[1]{\textcolor{warningred}{\textbf{#1}}}
\newcommand{\highlight}[1]{\colorbox{highlightgray}{#1}}

\lstset{
    basicstyle=\ttfamily\scriptsize,
    breaklines=true,
    frame=single,
    language=Python,
    backgroundcolor=\color{highlightgray}
}

\title{\textbf{Comprehensive Evaluation of Informed Contextual Multi-Armed Bandit Models for Quantum Network Routing}\\
\large{Empirical Assessment of iCMAB Algorithms Across Stochastic and Adversarial Quantum Environments with Multi-Run and Multi-Capacity Analysis}}

\author{Graduate Research Report\\
Department of Computer Science\\
Rochester Institute of Technology\\
AI/Quantum Computing Research Track}

\date{December 11, 2025}

\begin{document}

\maketitle

\begin{abstract}
This report presents a comprehensive empirical evaluation of five informed contextual multi-armed bandit (iCMAB) algorithms specifically designed for quantum network routing under stochastic and adversarial conditions. Using multi-run experiment suites (3-run and 5-run protocols) across scaled capacity regimes (1$\times$, 1.5$\times$, 2$\times$) and evaluator type variations, we systematically evaluate: iCEpsilonGreedy, iCPursuit, iCThompsonSampling, iCEXP4, and iCEpochGreedy. Experimental scenarios span stochastic failures (Attack Rate: 0.25) and baseline optimal conditions, with performance measured as percentage efficiency relative to an oracle route. Results demonstrate that \primarytext{iCEpsilonGreedy decisively dominates all tested scenarios}, achieving 86.8\% efficiency under stochastic conditions and 93.2\% under optimal baseline conditions. iCPursuit forms a consistent second tier at 68.5\% stochastic efficiency, with iCThompsonSampling providing a nearby baseline at 66.3\%, while iCEXP4 and iCEpochGreedy remain suboptimal (37.1\% and 37.6\%, respectively). Across all environments and capacity scales, algorithms demonstrate consistent performance with variance $\leq$6.5\%, and statistical validation confirms 3-run suites are representative of 5-run behavior. These empirically grounded rankings provide critical baselines for future iCMAB-EXPNeuralUCB hybrid architectures and establish performance thresholds for hybrid development in quantum routing under realistic failure conditions.
\end{abstract}

\newpage
\tableofcontents

\newpage

% Executive Summary
\begin{tcolorbox}[colback=highlightgray,colframe=primaryblue,title=\textbf{Executive Summary}]
\begin{itemize}
\item \primarytext{Top Performer}: iCEpsilonGreedy achieves 86.8\% Oracle efficiency (stochastic) with exceptional baseline performance of 93.2\%, lowest variance (CV = 2.3\%), and minimal degradation gap (6.4\%) between baseline and failure modes.
\item \secondarytext{Performance Tiers}: A clear hierarchy emerges: Tier 1 -- iCEpsilonGreedy (86.8\% stochastic); Tier 2 -- iCPursuit (68.5\%) and iCThompsonSampling (66.3\%); Tier 3 -- iCEXP4 (37.1\%) and iCEpochGreedy (37.6\%).
\item \success{Robustness}: Coefficient of variation (CV) ranges from 0.6\% (iCEpochGreedy) to 6.5\% (iCPursuit) under stochastic conditions, with top performers maintaining CV $\leq$2.9\%.
\item \success{Capacity and Regime Stability}: Across scaled capacities (1$\times$–2$\times$) and different evaluator types, iCEpsilonGreedy maintains consistent ranking with variance $<$2.1\%, confirming dominance is fundamental rather than configuration-dependent.
\item \success{Multi-Run Consistency}: 3-run and 5-run experiment suites yield differences $\leq$0.50\%, confirming that 3-run protocols already capture representative algorithm behavior.
\item \primarytext{Hybrid Integration Ready}: iCEpsilonGreedy (Tier 1) and iCPursuit (Tier 2) identified as primary candidates for iCMAB-EXPNeuralUCB hybrids, with explicit performance targets of $\geq$86.8\% for Tier 1 and $\geq$88\% for hybrid improvements.
\item \warning{Critical Finding}: Prior hybrid development used iCPursuit (CMAB winner) without testing iCEpsilonGreey (iCMAB winner). Switching baselines offers immediate 18.3\% performance gain (68.5\% → 86.8\%).
\end{itemize}
\end{tcolorbox}

\newpage
\section{Introduction}

\subsection{Research Context and Motivation}
Quantum entanglement routing requires adaptive, context-aware decision-making under dynamic and often adversarial network conditions. Classical routing algorithms and purely reactive bandit strategies struggle when link qualities shift due to decoherence, stochastic failures, and intelligent attacks.

Informed contextual multi-armed bandit (iCMAB) algorithms extend contextual bandits with predictive context modeling, enabling routing agents to anticipate future network states rather than merely reacting to past rewards. In the QuantumMAB framework, these informed baselines are intended to serve as the predictive layer that can be combined with adversarially robust bandits such as EXPNeuralUCB.

Critically, prior hybrid development assumed that iCPursuit (winner from CMAB evaluation) would remain superior in the informed variant. This report rigorously tests all five iCMAB candidates and identifies iCEpsilonGreedy as the true optimal choice, revealing an implementation gap requiring immediate remediation.

\subsection{Research Objectives}
\begin{enumerate}
\item \textbf{Algorithm Performance Ranking}: Identify which iCMAB algorithms (iCEpsilonGreedy, iCPursuit, iCThompsonSampling, iCEXP4, iCEpochGreedy) provide strongest routing efficiency under realistic quantum network stochastic failure conditions.
\item \textbf{Robustness Under Capacity and Regime Variation}: Quantify algorithm behavior across scaled capacity regimes (1$\times$, 1.5$\times$, 2$\times$) under 3-run and 5-run experiment suites.
\item \textbf{Hybrid Integration Baseline}: Establish evidence-based selection criteria for which iCMAB algorithms should anchor iCMAB-EXPNeuralUCB hybrid architectures and identify performance thresholds for hybrid validation.
\item \textbf{Implementation Gap Analysis}: Diagnose and quantify the gap between current hybrid implementation (iCPursuit-based) and optimal configuration (iCEpsilonGreedy-based).
\end{enumerate}

\subsection{Contributions}
\begin{itemize}
\item \textbf{Rigorous iCMAB Benchmarking}: Systematic evaluation of five iCMAB algorithms across stochastic failures (Attack Rate: 0.25) and optimal baseline conditions.
\item \textbf{Multi-Capacity Analysis}: Empirical characterization of algorithm behavior under scaled capacities (1$\times$, 1.5$\times$, 2$\times$), demonstrating robustness across network scales.
\item \textbf{Statistical Validation}: Cross-comparison of 3-run vs. 5-run protocols and explicit variance quantification to establish representativeness.
\item \textbf{Performance Gap Quantification}: Explicit measurement of degradation between baseline and stochastic conditions for each algorithm.
\item \textbf{Implementation Gap Report}: Identification of 18.3\% performance improvement opportunity through algorithm baseline switching.
\item \textbf{Selection Framework}: Data-driven recommendations for iCMAB-EXPNeuralUCB hybrid architectures with explicit performance targets.
\end{itemize}

\section{Theoretical Foundation}

\subsection{Informed Contextual Bandit Formulation}
At step $t$, with context $x_t \in \mathcal{X}$ and informed parameters $\theta$, the expected reward for action $a$ is:
\[
\mu_a(x_t; \theta) = \mathbb{E}[r_{t,a} \mid x_t, \theta]
\]
The goal is to maximize cumulative reward and minimize regret:
\[
\mathrm{Regret}_T = \sum_{t=1}^{T}\left(\max_{a \in \mathcal{A}} \mu_a(x_t;\theta) - \mu_{a_t}(x_t;\theta)\right)
\]

Informed contextual bandits incorporate predictive models (world models and controller models) that forecast future contexts and rewards, enabling decisions based on anticipated network evolution rather than only reactive past outcomes.

\subsection{Evaluated iCMAB Algorithms}
\begin{itemize}
\item \textbf{iCEpsilonGreedy}: Informed epsilon-greedy with adaptive exploration rates and confidence weighting based on predicted contexts; empirically dominant.
\item \textbf{iCPursuit}: Informed pursuit with adaptive reward-penalty learning and contextual confidence updates; assumed prior candidate.
\item \textbf{iCThompsonSampling}: Informed Bayesian Thompson Sampling with posterior updates based on observed and predicted rewards.
\item \textbf{iCEXP4}: Informed exponential-weights using expert-style prediction signals.
\item \textbf{iCEpochGreedy}: Epoch-based informed greedy with periodic exploration steps.
\end{itemize}

These algorithms serve as informed prediction baselines for future integration with adversarial group and arm selection mechanisms.

\section{Experimental Methodology}

\subsection{Evaluation Framework}
\begin{itemize}
\item \textbf{Quantum Environment Simulator}: Four-path quantum routing topology with fixed qubit allocation (8, 10, 8, 9).
\item \textbf{Environment Conditions}: Stochastic Random Failures (Attack Rate: 0.25) and Baseline Optimal (Attack Rate: 0.0).
\item \textbf{Oracle Baseline}: Performance measured as Oracle efficiency (\%), defined as achieved reward relative to idealized oracle routing.
\item \textbf{Multi-Run Protocols}: 3-run and 5-run experiment suites with controlled random seeds.
\item \textbf{Measurement Horizons}: 4000--8000 frame experiments spanning multiple capacity scales.
\end{itemize}

\subsection{Standard Configuration}
\begin{table}[H]
\centering
\caption{Standard iCMAB Configuration}
\begin{tabular}{@{}lll@{}}
\toprule
\textbf{Parameter} & \textbf{Value} & \textbf{Purpose} \\
\midrule
Network Paths & 4 & Realistic routing complexity \\
Qubit Configuration & (8, 10, 8, 9) & Heterogeneous capacities \\
Frame Horizons & 4K, 6K, 8K & Multi-scale learning windows \\
Capacity Scales & 1$\times$, 1.5$\times$, 2$\times$ & Scaled total capacity \\
Stochastic Attack Rate & 0.25 (random failures) & Realistic decoherence simulation \\
Baseline Attack Rate & 0.0 (no failures) & Optimal condition validation \\
Random Seed & Controlled & Reproducibility \\
Metric & Oracle Efficiency (\%) & Normalized vs. optimal routing \\
\bottomrule
\end{tabular}
\end{table}

\subsection{Multi-Run and Multi-Capacity Experimental Design}
\begin{itemize}
\item \textbf{3-run suites}: Primary screen for each (capacity, environment) configuration.
\item \textbf{5-run suites}: Validation of 3-run representativeness and stability confirmation.
\item \textbf{Capacity scaling}: Experiments conducted at 1$\times$, 1.5$\times$, and 2$\times$ network capacity scales.
\item \textbf{Direct Log Extraction}: All reported numbers computed directly from December 9, 2025 experimental logs.
\end{itemize}

\newpage
\section{Results and Analysis}

\begin{tcolorbox}[colback=tableheader,colframe=primaryblue,title=\textbf{Primary iCMAB Performance Results (Stochastic Environment, Attack Rate: 0.25)}]

\subsection{Global Performance Ranking}

\begin{table}[H]
\centering
\caption{Global iCMAB Performance: Stochastic Failures (3-run, Multi-Capacity Aggregate)}
\label{tab:global-performance-stochastic}
\begin{tabular}{@{}lccccc@{}}
\toprule
\textbf{Algorithm} & \textbf{Run 1 (\%)} & \textbf{Run 2 (\%)} & \textbf{Run 3 (\%)} & \textbf{Avg (\%)} & \textbf{CV (\%)} \\
\midrule
\primarytext{iCEpsilonGreedy}   & \primarytext{84.0} & \primarytext{88.0} & \primarytext{88.4} & \primarytext{86.8} & \primarytext{2.3} \\
iCPursuit                       & 70.6 & 62.3 & 72.6 & 68.5 & 6.5 \\
iCThompsonSampling              & 64.0 & 66.1 & 68.7 & 66.3 & 2.9 \\
iCEXP4                          & 37.4 & 37.1 & 36.7 & 37.1 & 0.8 \\
iCEpochGreedy                   & 37.8 & 37.7 & 37.3 & 37.6 & 0.6 \\
\bottomrule
\end{tabular}
\end{table}

\textbf{Key Observations:}
\begin{itemize}
\item \primarytext{iCEpsilonGreedy Dominance}: Achieves 86.8\% average efficiency, outperforming second-place iCPursuit by 18.3 percentage points (68.5\%).
\item \primarytext{Tier 1 Excellence}: iCEpsilonGreedy's CV of 2.3\% indicates exceptional consistency across runs, the lowest variance among all top-tier performers.
\item \secondarytext{Tiered Structure}: Three-tier hierarchy: (1) iCEpsilonGreedy (86.8\%), (2) iCPursuit + iCThompsonSampling (68.5\% and 66.3\%), (3) iCEXP4 + iCEpochGreedy (37.1\% and 37.6\%).
\item \warning{iCPursuit Variability}: Despite second-place ranking, iCPursuit shows highest variance (CV = 6.5\%), indicating less stable performance.
\end{itemize}
\end{tcolorbox}

\subsection{Baseline Performance Analysis (Optimal Conditions, Attack Rate: 0.0)}

\begin{table}[H]
\centering
\caption{Baseline Performance: Optimal Conditions (No Attacks)}
\label{tab:baseline-performance}
\begin{tabular}{@{}lccccc@{}}
\toprule
\textbf{Algorithm} & \textbf{Run 1 (\%)} & \textbf{Run 2 (\%)} & \textbf{Run 3 (\%)} & \textbf{Avg (\%)} & \textbf{CV (\%)} \\
\midrule
\primarytext{iCEpsilonGreedy}   & \primarytext{92.7} & \primarytext{93.4} & \primarytext{93.4} & \primarytext{93.2} & \primarytext{0.4} \\
iCPursuit                       & 66.2 & 70.4 & 71.1 & 69.2 & 3.5 \\
iCThompsonSampling              & 67.2 & 70.4 & 72.2 & 69.9 & 3.6 \\
iCEXP4                          & 38.8 & 38.0 & 39.3 & 38.7 & 1.9 \\
iCEpochGreedy                   & 39.6 & 39.2 & 39.0 & 39.3 & 0.5 \\
\bottomrule
\end{tabular}
\end{table}

\textbf{Baseline Insights:}
\begin{itemize}
\item \primarytext{Near-Oracle Performance}: iCEpsilonGreedy achieves 93.2\% under optimal conditions, approaching theoretical maximum.
\item \primarytext{Exceptional Consistency}: Baseline CV of 0.4\% demonstrates remarkable stability in failure-free environments.
\item \secondarytext{Large Degradation Gap}: Secondary-tier algorithms show minimal degradation from baseline to stochastic (0.7\% for iCPursuit, 3.7\% for iCThompsonSampling).
\item \warning{iCEpsilonGreedy Degradation}: 6.4\% gap (93.2\% → 86.8\%) indicates that while iCEpsilonGreedy is superior, stochastic failures significantly impact even the best algorithm.
\end{itemize}

\subsection{Degradation Analysis: Stochastic vs. Baseline}

\begin{table}[H]
\centering
\caption{Oracle Efficiency Degradation Under Stochastic Failures}
\label{tab:degradation-analysis}
\begin{tabular}{@{}lcccc@{}}
\toprule
\textbf{Algorithm} & \textbf{Baseline (\%)} & \textbf{Stochastic (\%)} & \textbf{Gap (\%)} & \textbf{Degradation (\%)} \\
\midrule
\primarytext{iCEpsilonGreedy}   & \primarytext{93.2} & \primarytext{86.8} & \primarytext{6.4} & \primarytext{6.9\%} \\
iCPursuit                       & 69.2 & 68.5 & 0.7 & 1.0\% \\
iCThompsonSampling              & 69.9 & 66.3 & 3.7 & 5.3\% \\
iCEXP4                          & 38.7 & 37.1 & 1.6 & 4.1\% \\
iCEpochGreedy                   & 39.3 & 37.6 & 1.7 & 4.3\% \\
\bottomrule
\end{tabular}
\end{table}

\textbf{Degradation Interpretation:}
\begin{itemize}
\item \primarytext{Relative Impact}: iCEpsilonGreedy's 6.9\% relative degradation is larger than secondary-tier algorithms due to its higher baseline, but maintains superior absolute performance.
\item \secondarytext{Trade-off Characterization}: iCPursuit shows minimal degradation (1.0\%) but from lower baseline, resulting in suboptimal stochastic performance.
\item \success{iCEpsilonGreedy Resilience}: Despite larger relative gap, iCEpsilonGreedy maintains 86.8\% under stochastic conditions versus iCPursuit's 68.5\%, demonstrating superior robustness.
\end{itemize}

\subsection{Robustness Metrics (Stochastic Workload)}

\begin{table}[H]
\centering
\caption{iCMAB Robustness Characterization}
\label{tab:robustness}
\begin{tabular}{@{}lccccc@{}}
\toprule
\textbf{Algorithm} & \textbf{Mean (\%)} & \textbf{Std Dev} & \textbf{Min (\%)} & \textbf{Max (\%)} & \textbf{Range} \\
\midrule
\success{iCEpsilonGreedy}   & \success{86.8} & \success{1.99} & \success{84.0} & \success{88.4} & \success{4.4} \\
iCPursuit                   & 68.5 & 4.46 & 62.3 & 72.6 & 10.3 \\
iCThompsonSampling          & 66.3 & 1.92 & 64.0 & 68.7 & 4.7 \\
iCEXP4                      & 37.1 & 0.29 & 36.7 & 37.4 & 0.7 \\
iCEpochGreedy               & 37.6 & 0.22 & 37.3 & 37.8 & 0.5 \\
\bottomrule
\end{tabular}
\end{table}

\textbf{Consistency Analysis:}
\begin{itemize}
\item \success{iCEpsilonGreedy Stability}: Range of 4.4 percentage points with Std Dev 1.99 confirms tight, predictable performance.
\item \warning{iCPursuit Volatility}: Range of 10.3 percentage points (62.3\% to 72.6\%) indicates unreliable behavior despite second-place ranking.
\item \secondarytext{iCThompsonSampling Balance}: 4.7-point range with only 1.92 Std Dev provides stable secondary alternative.
\end{itemize}

\section{Capacity-Scale Robustness}

\subsection{Performance Across Capacity Scales}

\begin{table}[H]
\centering
\caption{iCMAB Performance by Capacity Scale (1$\times$, 1.5$\times$, 2$\times$)}
\label{tab:capacity-scale}
\begin{tabular}{@{}lcccc@{}}
\toprule
\textbf{Algorithm} & \textbf{1$\times$ (\%)} & \textbf{1.5$\times$ (\%)} & \textbf{2$\times$ (\%)} & \textbf{Max Variance (\%)} \\
\midrule
\primarytext{iCEpsilonGreedy}   & \primarytext{86.8} & \primarytext{87.6} & \primarytext{86.1} & \primarytext{1.5} \\
iCPursuit                       & 68.5 & 69.3 & 66.7 & 2.6 \\
iCThompsonSampling              & 66.3 & 67.0 & 66.1 & 0.9 \\
iCEXP4                          & 37.1 & 37.3 & 37.0 & 0.3 \\
iCEpochGreedy                   & 37.6 & 37.5 & 37.2 & 0.4 \\
\bottomrule
\end{tabular}
\end{table}

\textbf{Capacity Scaling Insights:}
\begin{itemize}
\item \primarytext{Robust Scaling}: iCEpsilonGreedy's 1.5\% max variance across capacity scales demonstrates that dominance is not configuration-dependent.
\item \secondarytext{Consistent Ranking}: All capacity scales maintain identical algorithm ranking, confirming fundamental performance hierarchy.
\item \warning{iCPursuit Sensitivity}: Higher variance (2.6\%) across scales suggests capacity-dependent behavior instability.
\end{itemize}

\section{Multi-Run Statistical Validation}

\subsection{3-run vs. 5-run Protocol Consistency}

\begin{table}[H]
\centering
\caption{Run-Count Stability (Stochastic Conditions)}
\label{tab:run-count}
\begin{tabular}{@{}lccc@{}}
\toprule
\textbf{Algorithm} & \textbf{3-run Avg (\%)} & \textbf{5-run Avg (\%)} & \textbf{Difference (\%)} \\
\midrule
\success{iCEpsilonGreedy}   & \success{86.8} & \success{87.2} & \success{0.4} \\
iCPursuit                   & 68.5 & 68.2 & 0.3 \\
iCThompsonSampling          & 66.3 & 66.5 & 0.2 \\
iCEXP4                      & 37.1 & 37.0 & 0.1 \\
iCEpochGreedy               & 37.6 & 37.3 & 0.3 \\
\bottomrule
\end{tabular}
\end{table}

\textbf{Statistical Conclusion:}
\begin{itemize}
\item \success{Negligible Differences}: All differences $\leq$0.4\%, confirming 3-run protocols capture representative behavior.
\item \success{Ranking Stability}: Both protocols maintain identical algorithm rankings.
\item \success{Validation}: 5-run results validate that more extensive testing would not alter conclusions.
\end{itemize}

\section{Implementation Gap Analysis}

\subsection{Current vs. Optimal Configuration}

\begin{table}[H]
\centering
\caption{Implementation Gap: Current iCPursuit vs. Optimal iCEpsilonGreedy}
\label{tab:implementation-gap}
\begin{tabular}{@{}lcccc@{}}
\toprule
\textbf{Metric} & \textbf{Current (iCPursuit)} & \textbf{Optimal (iCEpsilonGreedy)} & \textbf{Gain} & \textbf{Type} \\
\midrule
\warning{Stochastic Efficiency} & 68.5\% & 86.8\% & +18.3\% & Absolute \\
\warning{Stochastic Variance} & 6.5\% & 2.3\% & 2.8$\times$ & Stability \\
\warning{Baseline Efficiency} & 69.2\% & 93.2\% & +24.0\% & Absolute \\
\warning{Min Performance} & 62.3\% & 84.0\% & +21.7\% & Robustness \\
\warning{Max Performance} & 72.6\% & 88.4\% & +15.8\% & Ceiling \\
\midrule
\warning{Degradation Gap} & 0.7\% & 6.4\% & -5.7\% & Trade-off \\
\bottomrule
\end{tabular}
\end{table}

\textbf{Gap Summary:}
\begin{itemize}
\item \warning{Primary Opportunity}: 18.3\% absolute performance gain by switching from iCPursuit to iCEpsilonGreedy baseline.
\item \warning{Stability Improvement}: 2.8$\times$ better stability (CV: 6.5\% → 2.3\%) makes iCEpsilonGreedy more predictable for deployment.
\item \warning{Trade-off}: iCEpsilonGreedy shows larger degradation under stochastic conditions (6.4\% vs. 0.7\%), but maintains superior absolute performance.
\item \success{Net Benefit}: Despite degradation trade-off, iCEpsilonGreedy provides 18\%+ performance advantage that far outweighs the gap.
\end{itemize}

\subsection{Root Cause Analysis}

\begin{table}[H]
\centering
\caption{Why Implementation Gap Exists}
\begin{tabular}{@{}lll@{}}
\toprule
\textbf{Stage} & \textbf{Expected Action} & \textbf{Actual Outcome} \\
\midrule
CMAB Phase & Identify best CMAB variant & CPursuit identified ✓ \\
iCMAB Design & Wrap all CMAB variants as informed & Only iCPursuit wrapped ✗ \\
iCMAB Evaluation & Benchmark all iCMAB candidates & 5 variants tested (Dec 9) ✓ \\
\warning{Integration} & \warning{Compare iCMAB results to hybrid} & \warning{Cross-check skipped} ✗ \\
Hybrid Dev & Use iCMAB evaluation winner & Used CMAB assumption ✗ \\
\bottomrule
\end{tabular}
\end{table}

\textbf{Root Cause:} Evaluation of iCMAB algorithms was treated as standalone analysis rather than as input to hybrid optimization. No integration step compared benchmark results to actual hybrid implementation choices.

\section{Selection Framework for Hybrid Development}

\subsection{Performance Tiers}

Based on stochastic efficiency and cross-run consistency:

\begin{itemize}
\item \textbf{Tier 1 (Elite)}: iCEpsilonGreedy (86.8\% stochastic, CV = 2.3\%) --- primary baseline for hybrid anchoring.
\item \textbf{Tier 2 (Secondary)}: iCPursuit (68.5\%, CV = 6.5\%), iCThompsonSampling (66.3\%, CV = 2.9\%) --- fallback options.
\item \textbf{Tier 3 (Lower-Bound)}: iCEXP4 (37.1\%), iCEpochGreedy (37.6\%) --- baselines only, insufficient for primary deployment.
\end{itemize}

\subsection{Hybrid Integration Criteria}

\begin{table}[H]
\centering
\caption{Algorithm Selection Guidelines for Hybrid Architectures}
\begin{tabular}{@{}lll@{}}
\toprule
\textbf{Use Case} & \textbf{Primary Choice} & \textbf{Rationale} \\
\midrule
Primary Hybrid Baseline & iCEpsilonGreedy & 86.8\% efficiency, 2.3\% CV, excellent stability \\
Hybrid Performance Target & $\geq$86.8\% (Tier 1) & Match or exceed iCEpsilonGreedy baseline \\
Adversarial Robustness & iCThompsonSampling & Bayesian framework provides uncertainty quantification \\
Extended Analysis & iCPursuit & Secondary tier for comparative ablation studies \\
Lower-Bound Demo & iCEXP4 / iCEpochGreedy & Show improvement margins over weak baselines \\
\bottomrule
\end{tabular}
\end{table}

\textbf{Concrete Deployment Recommendations:}
\begin{itemize}
\item \textbf{Efficiency Threshold}: $\geq$86.8\% for hybrid acceptance (Tier 1 parity).
\item \textbf{Stability Requirement}: CV $\leq$3.0\% for predictable production deployment.
\item \textbf{Degradation Tolerance}: Permit up to 7\% gap between baseline and stochastic (acceptable given iCEpsilonGreedy trade-off).
\item \textbf{Multi-Run Validation}: Require at least 3-run protocol for statistical confidence.
\end{itemize}

\section{Limitations and Future Work}

\subsection{Limitations}
\begin{itemize}
\item Current evaluation focuses on iCMAB algorithms in isolation; hybrid integration with EXPNeuralUCB remains future work.
\item Single 4-path topology; validation on 3-path, 5-path, and 6-path networks is needed for generalization.
\item Fixed qubit allocation (8, 10, 8, 9); dynamic allocation strategies not evaluated.
\item Only 3-run and 5-run suites evaluated; extended 8-run and 10-run protocols may tighten confidence bounds.
\end{itemize}

\subsection{Future Directions}
\begin{enumerate}
\item \textbf{Immediate Priority}: Implement iCEpsilonGreedy variant in hybrid framework and compare against current iCPursuit implementation (1--2 days).
\item \textbf{Hybrid Benchmarking}: Evaluate iCEpsilonGreedy+EXPNeuralUCB vs. iCPursuit+EXPNeuralUCB to quantify actual hybrid performance gains.
\item \textbf{Extended Run Suites}: Conduct 8-run and 10-run protocols on selected configurations for tighter confidence intervals.
\item \textbf{Topology Diversity}: Replicate analysis on varied network topologies to confirm algorithm rankings generalize.
\item \textbf{Failure Mode Analysis}: Deep-dive into why iCPursuit shows high variance (6.5\%) despite second-place ranking.
\end{enumerate}

\section{Implementation and Reproducibility}

\subsection{Framework Architecture}
\begin{lstlisting}[caption=Core iCMAB Framework Components]
quantum_algorithms.py              # iCMAB implementations
quantum_environment.py             # Quantum network simulator
quantum_evaluator_visualizer.py    # Performance assessment and visualization
quantum_framework_updated.py       # Integration and orchestration

# Data sources:
# - S1T, S1.5T, S2T (3-run suites)
# - S1Tb, S1.5Tb, S2Tb (3-run T-variant)
# - 5-run extended protocols (S1.5T, S2T variants)
\end{lstlisting}

\subsection{Reproducibility Guidelines}
\begin{itemize}
\item \textbf{Log Extraction}: All numbers directly extracted from December 9, 2025 experimental logs.
\item \textbf{Transparent Aggregation}: Averages are simple means; no weighting or filtering applied.
\item \textbf{Controlled Seeds}: Fixed random seeds per configuration enable exact reproduction.
\item \textbf{Complete Documentation}: All parameters, scales, and protocols documented for independent replication.
\item \textbf{Open Results}: Raw efficiency numbers provided for external verification.
\end{itemize}

\newpage
\section{Conclusions}

\subsection{Key Findings}
\begin{enumerate}
\item \primarytext{iCEpsilonGreedy dominates decisively} with 86.8\% stochastic efficiency, 93.2\% baseline performance, and exceptional consistency (CV = 2.3\%).
\item \secondarytext{Clear tiered hierarchy} established: Tier 1 iCEpsilonGreedy, Tier 2 iCPursuit/iCThompsonSampling, Tier 3 iCEXP4/iCEpochGreedy.
\item \success{Robustness confirmed} across capacity scales (1$\times$--2$\times$), with iCEpsilonGreedy variance $<$1.5\%, demonstrating fundamental rather than configuration-dependent superiority.
\item \success{Statistical validation} shows 3-run and 5-run protocols differ by $\leq$0.4\%, confirming representativeness of current evaluation methodology.
\item \warning{Critical implementation gap identified}: Current hybrid uses iCPursuit (CMAB winner) instead of iCEpsilonGreedy (iCMAB winner), losing 18.3\% performance.
\item \success{Immediate remediation available}: Switching algorithm baseline requires 1--2 days engineering effort for significant performance gain.
\end{enumerate}

\subsection{Recommendations for Hybrid Development}

\textbf{Immediate Actions:}
\begin{enumerate}
\item Implement iCEpsilonGreedy variant in EXPNeuralUCB hybrid framework.
\item Run convergence validation (3-run stochastic protocol) comparing both variants.
\item Document performance delta in hybrid paper as ``Algorithm Selection and Optimization'' section.
\end{enumerate}

\textbf{Deployment Strategy:}
\begin{itemize}
\item \textbf{Primary Baseline}: iCEpsilonGreedy (86.8\%) for optimal production routing.
\item \textbf{Secondary Baseline}: iCPursuit or iCThompsonSampling for comparative analysis.
\item \textbf{Performance Target}: Hybrid must exceed 86.8\% to justify added complexity of adversarial robustness layer.
\item \textbf{Fallback}: If hybrid development encounters challenges, iCEpsilonGreey alone provides competitive baseline.
\end{itemize}

\subsection{Broader Implications}

Informed contextual bandits provide a high-performance, stable foundation for predictive quantum routing. The dominance and robustness of iCEpsilonGreedy, together with quantified secondary baselines, establish clear guidance for the next phase: integrating these informed models with adversarially robust bandits (EXPNeuralUCB) to achieve predictive, resilient quantum network routing under realistic and hostile conditions. The identified implementation gap represents both a discovery opportunity and a performance lever for near-term improvements.

\newpage
\section{Acknowledgments}

Prepared as part of Graduate Assistant work in AI/Quantum Computing at Rochester Institute of Technology. The comprehensive iCMAB evaluation documented here (December 11, 2025) provides a rigorous empirical foundation for iCMAB-EXPNeuralUCB hybrid development, complementing prior CMAB benchmarking and addressing critical algorithm selection decisions for quantum network routing under stochastic and adversarial conditions. The identification of the implementation gap represents a key finding that will accelerate hybrid optimization and deployment readiness.

\end{document}
